\section{Methodology and Experimental Setup}
\label{sec:method}
\subsection{Platforms, datasets and correctness checks}
The primary dataset contains CPU library measurements from two environments. The Windows system used an Intel Core i5-10300H with four physical cores and eight logical processors, Python 3.12.10, PyCryptodome 3.23.0 and verified 64-bit OpenSSL 3.6.5 libraries. The Oracle Cloud system used a Neoverse-N1-class Arm64 processor with two available virtual CPUs, Python 3.12.3 and OpenSSL 3.0.13. The original records preserve platform metadata, processor capabilities and source hashes. The hosts differ in operating system, library version and hardware configuration, so each is analyzed separately.

The primary campaigns evaluated AES-128 at payload sizes of 1 KiB, 16 KiB, 1 MiB, 16 MiB and 32 MiB. CTR is the principal workload presented here. Electronic codebook (ECB) measurements in the source records served as implementation diagnostics. Each configuration was measured in five sessions with seeds 2003--2007, using five warm-up repetitions and 20 measured repetitions per session. Known-answer and differential checks established ciphertext correctness before performance results were accepted. Raw observations, metadata and session summaries were retained.

Two earlier Windows campaigns provide complementary evidence. The CBC experiment compared software and AES-NI-enabled implementations for 128-, 192- and 256-bit keys; the present paper reports its 10 MiB encryption results. The GCM experiment evaluated 128- and 256-bit keys at 1 KiB, 16 KiB, 1 MiB, 16 MiB and 100 MiB. Both used five seeded sessions, three warm-up repetitions and 20 measured repetitions. These datasets retain their original timing definitions and are analyzed separately from the CTR campaign.

The GCM workload used a 96-bit nonce, a 128-bit authentication tag and 16 bytes of associated authenticated data (AAD), with no nonce reuse for distinct messages under the same key within a session. Validation included known-answer checks and rejection of modified ciphertext, tags and AAD. These checks establish functional correctness for the tested cases; they do not constitute a side-channel security evaluation.

\subsection{Implementation paths and timing scopes}
The CTR comparison uses PyCryptodome and native OpenSSL EVP. Both transform measurements use prepared output buffers. For OpenSSL, timestamps inside the compiled wrapper delimit the EVP Update/Final processing. For PyCryptodome, the transform interval includes Python dispatch and the library's internal foreign-function interface. Thus, the recorded transform rates describe each implementation at its measured interface; their ratio does not isolate the relative cost of AES instructions.

A second interval, termed the \emph{instrumented operation}, measures the Python call to the benchmark's measurement routine. It includes input checks, allocation, buffer bindings, setup, transformation, cleanup and output conversion, together with diagnostic calls and timing bookkeeping where applicable. This is the interval named \texttt{operation\_envelope} in the records. It provides a broader view of the same benchmark execution, but is not an uninstrumented application request or a complete storage or network operation. The Windows dataset provides both intervals for the library comparison.

On Arm64, OpenSSL was also evaluated with \texttt{OPENSSL\_armcap=0}. This setting disables multiple optimized capability paths. Comparing it with the default configuration tests sensitivity to capability selection within OpenSSL, but cannot isolate the effect of the AES extension from other disabled features. The analysis uses the archived Arm64 transform summaries and does not infer unrecorded instruction execution from capability flags.

The CBC results compare the software and AES-NI paths under the original campaign's common timing convention. GCM encryption includes AAD processing, payload encryption, output allocation and tag generation; decryption includes authentication-tag verification. Construction of a fresh GCM cipher object is excluded. Throughput counts payload bytes only, even though AAD and tag processing contribute to elapsed time. These rates therefore describe authenticated payload processing after object construction, rather than complete message latency.

\subsection{Aggregation and interpretation}
For a payload of $B$ bytes processed in $t$ nanoseconds, decimal throughput is
\begin{equation}
 P=\frac{10^3 B}{t}\quad [\mathrm{MB/s}].\label{eq:rate}
\end{equation}
KiB and MiB denote binary payload sizes; MB/s and GB/s denote decimal rates. Each primary result is the median of five session-level throughput medians. Session ranges, where reported, describe the smallest and largest of those five medians rather than confidence intervals.

For a reference implementation $A$ and a candidate $B$ measured under a common timing convention, the speedup of $B$ relative to $A$ is
\begin{equation}
 S_{B/A}^{(j)}=\frac{\widetilde t_A^{(j)}}{\widetilde t_B^{(j)}},\qquad
 S_{B/A}=\operatorname{median}_{j=1}^{5}S_{B/A}^{(j)}.\label{eq:ratio}
\end{equation}
Here, $\widetilde t$ denotes the median elapsed time in session $j$; a value above one indicates faster execution by $B$. The CBC acceleration factors use this estimator. Approximate CTR throughput ratios discussed below are ratios of the reported aggregate rates, and are distinguished from session-based speedups. Rounding and the order of aggregation can make these estimators differ.

The session is the unit of replication. Repetitions are not pooled to imply additional independent experiments, and observations are not excluded because they show unfavorable performance. Five sessions characterize repeatability on the measured host, but do not establish population-level behavior across processors. Comparisons across hosts or cipher modes are descriptive and are not treated as paired trials.
